\section{Introduction: models of the soil water balance}
Chapter~3 gave the physics of water flow in a soil: the Richards equation. For many purposes (scheduling, design, planning) that is more detail than needed, and a simpler model is used: the soil is a \emph{bucket} whose content changes each day by what comes in and what goes out. Models of the soil water balance can be ranked by their detail. The \emph{tipping bucket} has one store, refilled by rain and irrigation and emptied by evapotranspiration, and spills the excess as drainage. The \emph{cascading bucket} has several stores in series, each spilling to the next. The \emph{capacity models} add a rate of drainage that depends on the content. The \emph{mechanistic} models solve Richards' equation. The application uses two engines: Engine~1, which reproduces the relations of the Daisy ensemble by ratios, and Engine~2, a daily bucket with the FAO stress function; the dashboard presents their results as indicators.

\section{Assumptions and equations of the tipping bucket}
The root zone is one layer of depth $Z_r$ with capacity $TAW$ between field capacity and wilting point. Each day $i$: (1) rain $P_i$ is reduced by interception and runoff to the infiltration $f_i$; (2) irrigation $I_i$ is added; (3) the actual evapotranspiration is $ET_{a,i}=K_{s,i}K_{c,i}ET_{0,i}$; (4) the depletion is updated and the excess over field capacity drains in the same day. The depletion is

\begin{equation}D_i=\min\bigl\{TAW,\ \max\{0,\ D_{i-1}-f_i-I_i+ET_{a,i}\}\bigr\},\qquad DP_i=\max\{0,\ f_i+I_i-ET_{a,i}-D_{i-1}\},\label{eq:bucket8}\end{equation}


with the stress coefficient $K_s=1$ for $D\le RAW$ and $K_s=(TAW-D)/(TAW-RAW)$ beyond it. The model assumes that drainage is instantaneous, that there is no capillary rise and that the root zone is homogeneous; these hold best for deep, freely drained soils.

\section{Runoff: the curve-number method}
Rain that arrives faster than the soil can take it, or on a saturated soil, runs off. The \kw{curve-number} method of the US Soil Conservation Service estimates the direct runoff $Q$ from a storm of depth $P$ by

\begin{equation}Q=\frac{(P-I_a)^2}{P-I_a+S},\quad I_a=0.2S,\quad S=\frac{25\,400}{CN}-254\ \ (\text{mm}),\ \ P>I_a,\label{eq:cn8}\end{equation}


with the curve number $CN$ from 30 (very permeable, dry) to 100 (impervious). A cultivated loam in fair condition has $CN\approx75$ at average antecedent moisture, 87 when the soil is wet and 57 when dry.

\exbox{Runoff of a storm}{A storm of 60 mm falls on a loam with $CN=75$. Find $S$, $I_a$ and the runoff. How does the runoff change if the soil is wet ($CN=87$)?}\sol \givenline{$CN=75$ ; $CN=87$}\ans $S=25400/75-254=84.7$ mm; $I_a=0.2S=16.9$ mm; $Q=(60-16.9)^2/(60-16.9+84.7)=14.5$ mm \ans For $CN=87$: $S=38.0$, $I_a=7.6$, $Q=30.4$ mm \ans Wet soil converts 51 per cent of the storm to runoff, against 24 per cent when it is average. \finans

\section{A season of the bucket, by program}
The bucket is easy to program, and a season can be run in a fraction of a second. The program of this book (and the engine of the application) reads a daily record of $ET_0$ and rain, builds $K_c(t)$ from the four stages, applies the irrigation rule (a trigger on depletion with a dose that refills to field capacity) and writes the totals: $ET_a$, irrigation, drainage, runoff and the relative yield from the Doorenbos--Kassam relation. The result of one season is a number of tables and charts; the design question is what the best trigger is. Section~\ref{sec:proj8} below compares triggers for a maize season.

\section{The dashboard of the application}
The dashboard is the summary page of the application: one set of indicators read from the layers. The most important are (i) the \emph{irrigation} applied, (ii) the \emph{evapotranspiration} and its split into $T$, $E_s$ and $E_o$, (iii) the \emph{drainage} and the nitrogen it carries, (iv) the \emph{yield} and the \emph{water productivity} $WP=Y/ET_a$, (v) the \emph{irrigation efficiency} $IE=(ET_a-ET_{rain})/I$ (the share of the applied water that the crop transpired), and (vi) the \emph{closure} of the balance. Each indicator has a formula, and the chapter on the dashboard shows how every number is produced; here we compute two of them by hand.

\exbox{Dashboard indicators}{A season gives $P=380$, $I=240$, $ET_a=520$, $D=70$, $Q=10$ mm, yield 11.2 Mg\,ha$^{-1}$, and the rainfed twin has $ET_a=410$. Find WP, the irrigation efficiency based on the extra evapotranspiration, and the closure.}\sol \givenline{$P=380$ ; $I=240$ ; $ET_a=520$ ; $D=70$ ; $Q=10$ mm ; $ET_a=410$}\ans $WP=11\,200/5\,200=2.15$ kg\,m$^{-3}$ \ans $IE=(520-410)/240=0.458$ \ans Balance: $P+I=620$; outputs $ET+D+Q=600$, so $\Delta S=20$ mm if the storage has changed by that amount; the residual is zero by definition once $\Delta S$ is entered \ans Only 46 per cent of the irrigation became extra evapotranspiration; the rest drained, ran off or stayed in the soil. \finans

